\documentclass[10pt,a4paper]{amsart}
\usepackage[margin=19mm]{geometry}
\usepackage{amsmath,amssymb,mathtools}
\usepackage[scr=rsfs,cal=euler]{mathalfa}
\usepackage{xfrac}
\usepackage[unicode,hidelinks,bookmarks=false]{hyperref}
\newcommand{\ie}{i.e.\ }
\newcommand{\cf}{cf.\ }
\newcommand{\resp}{respectively\ }
\newcommand{\Subsection}{\subsection}
\providecommand{\nobreakdash}{\nobreak\hbox{-}}
\setlength{\emergencystretch}{2em}
\multlinegap0pt
\usepackage{braket}
\usepackage{amssymb}
\usepackage{xcolor}
\usepackage{tikz}
\usepackage{algorithm}\usepackage{algpseudocode}
\usepackage[normalem]{ulem}
\newtheorem{thm}{}
\newtheorem{cor}{}
\newcommand{\Tr}{\mathrm{Tr}}
\title{Quantum Gibbs sampling through the detectability lemma}
\author{Di Fang, Jianfeng Lu, Yu Tong, Chu Zhao}
\date{Source: arXiv:2604.07214v1 (CC BY 4.0)}
\begin{document}
\maketitle

\noindent\emph{Source and licence.} Quantum Gibbs sampling through the detectability lemma. Version arXiv:2604.07214v1 (CC BY 4.0), distributed by arXiv under the Creative Commons Attribution 4.0 International licence, http://creativecommons.org/licenses/by/4.0/, as recorded in the arXiv licence metadata for this exact version and shown on the abstract page https://arxiv.org/abs/2604.07214v1. Full licence notice: \texttt{licenses/arxiv-2604.07214-CC-BY-4.0.txt}.\par\smallskip

\noindent\textit{Editorial scope.} Counted unit: the article's cor cor:prepare\_gibbs\_state\_quantum\_gibbs\_sampler (source0.tex lines 321--328). The article's complete original proof of it is delivered with it (source0.tex lines 330--348, one \texttt{proof} environment of 19 lines). No mathematics is written, repaired or re-derived: the statement and the proof are the article's own lines, in the article's own notation, and no retained line is substituted or re-flowed.  Retained uncounted as the source-local context the proof needs: lines 267--270; lines 311--318.  Nothing was omitted from any retained span, so the package carries no omission record.  No retained line contains a citation, so the wrapper carries no bibliography and no bibliography entry.  No retained line contains a figure environment, an image-inclusion command or drawing code, so no image is delivered and the package declares no asset dependency.  Editorial formatting only, in addition: an AMS \texttt{amsart} preamble that restates the article's own declarations and macros used by the retained lines, namely the environments \texttt{thm}, \texttt{cor} on separate counters, together with the standard packages those lines need.  The retained environments print in this document as Theorem 1, Corollary 1 in order of appearance, whereas the article prints the same items with the numbering of its own full document.  The complete native source of the article as distributed in the arXiv e-print for this exact version is bundled unchanged as \texttt{sources/198075/source0.tex}.
\begin{equation}
\label{eq:detectability_lem_op}
    \Phi = P_1 P_2\cdots P_M,\quad P_m = \lim_{t\to \infty} e^{t\mathcal{L}_m},\quad \forall 1\leq m\leq M.
\end{equation}

\begin{thm}
    \label{thm:quantum_gibbs_sampler_shrink}
    Let $\mathcal{L}=\sum_{m=1}^M \mathcal{L}_m$, where each $\mathcal{L}_m$ is a Lindbladian that satisfies $\sigma$-KMS DBC, and we assume $\mathcal{L}$ has a non-degenerate $0$-eigenspace. We assume that each $\mathcal{L}_m$ commutes with all but at most $g$ others. Then let the quantum channel $\Phi$ be defined as in \eqref{eq:detectability_lem_op}. For any $X\in B(\mathcal{H})$ such that $\Tr[\sigma X]=0$, we have
    \[
    \|\Phi(X)\|_\sigma^2\leq \frac{\|X\|_\sigma^2}{\mathrm{gap}(\mathcal{L})/g^2+1}.
    \]
    Moreover $\Tr[\sigma \Phi(X)]=\Tr[\Phi^\dag(\sigma) X]=\Tr[\sigma X]=0$.
\end{thm}

\begin{cor}
    \label{cor:prepare_gibbs_state_quantum_gibbs_sampler}
    Under the same assumptions as Theorem~\ref{thm:quantum_gibbs_sampler_shrink}, let $\rho_k=(\Phi^\dag)^k(\rho_0)$ for any initial state $\rho_0$, then
    \[
    \|\rho_k-\sigma\|_1\leq \frac{1}{(\mathrm{gap}(\mathcal{L})/g^2+1)^{k/2}\sqrt{\sigma_{\min}}},
    \]
    where $\sigma_{\min}$ is the smallest eigenvalue of $\sigma$.
\end{cor}

\begin{proof}
    Let $X\in B(\mathcal{H})$ satisfy $\|X\|\leq 1$. Denote $\Tilde{X} = X - \Tr[\sigma X] I$, and we have $\|\Tilde{X}\|_\sigma\leq \|X\|_\sigma\leq \|X\|\leq 1$. This ensures $\Tr[\sigma\tilde{X}]=0$. We then compute
    \[
    \Tr[X(\rho_k-\sigma)] = \Tr[\Tilde{X}(\rho_k-\sigma)]=\Tr[\Tilde{X}\rho_k]=\Tr[\Tilde{X}(\Phi^\dag)^k(\rho_0)] = \Tr[\Phi^k(\Tilde{X})\rho_0].
    \]
    By the Cauchy-Schwarz inequality, we have
    \[
    |\Tr[\Phi^k(\Tilde{X})\rho_0]|=|\braket{\Phi^k(\Tilde{X}),\sigma^{-1/2}\rho_0\sigma^{-1/2}}_\sigma|\leq \|\Phi^k(\Tilde{X})\|_\sigma\|\sigma^{-1/2}\rho_0\sigma^{-1/2}\|_\sigma.
    \]
    For $\|\Phi^k(\Tilde{X})\|_\sigma$ we directly apply Theorem~\ref{thm:quantum_gibbs_sampler_shrink} $k$ times and use the fact that $\|\Tilde{X}\|_\sigma\leq 1$ and $\Tr(\sigma \Tilde{X}) = 0$. For $\|\sigma^{-1/2}\rho_0\sigma^{-1/2}\|_\sigma$ we have
    \[
    \|\sigma^{-1/2}\rho_0\sigma^{-1/2}\|_\sigma = \Tr[\rho_0\sigma^{-1/2}\rho_0\sigma^{-1/2}]^{1/2}\leq \frac{1}{\sqrt{\sigma_{\min}}}.
    \]
    Therefore we have
    \[
    |\Tr[X(\rho_k-\sigma)]|\leq \frac{1}{(\mathrm{gap}(\mathcal{L})/g^2+1)^{k/2}\sqrt{\sigma_{\min}}}.
    \]
    Because this is true for any $\|X\|\leq 1$, we have the desired inequality for the trace distance between $\rho_k$ and $\sigma$ through the duality between the Schatten $1$- and $\infty$-norms.
\end{proof}

\end{document}
